\begin{figure}[ht]
  \centering
  \includegraphics[width=\textwidth]{figures/fig_pathway_delta_dotplot.pdf}
  \caption{\textbf{Rank-difference GSEA: which gene sets each deconfounding
  method promotes or demotes relative to marginal.} (\textbf{a}) DepMap
  (essentiality); (\textbf{b}) CTRPv2 (drug response). For each deconfounding
  method (linear residualization, IRM, within-tissue SHAP) we build a per-gene
  promotion score
  $\delta_g = \mathrm{r}_{\mathrm{method}}(g) - \mathrm{r}_{\mathrm{marginal}}(g)$,
  where $\mathrm{r}(g)$ is gene $g$'s within-item rank of $|\mathrm{SHAP}|$
  (top $=1$) divided by the panel size and averaged across items (drugs or
  knockout targets), so each item contributes equally regardless of its
  $|\mathrm{SHAP}|$ magnitude. A positive $\delta_g$ means $g$ is ranked higher
  under the method than under marginal SHAP. We then run pre-ranked GSEA
  (gseapy, 10{,}000 permutations) of the KEGG, Hallmark and Human Gene Atlas
  (HGA) libraries on the $\delta$ vector, excluding KEGG ``Human Diseases'' pathways other than cancer (for example infection and substance-dependence maps). Upper plot: gene sets reaching FDR $q<0.05$ in at least one of the
  three deconfounding methods for that task (DepMap $n=31$, CTRPv2 $n=37$),
  ordered by mean NES across the three methods. Dot color is the promotion NES
  (red, promoted; blue, demoted) and area is proportional to
  $-\log_{10}(\mathrm{FDR}\,q)$. The left strip and label color give the
  gene-set class, assigned from KEGG BRITE and MSigDB
  Hallmark process categories: orange, HGA tissue (lineage) signature; gold,
  tissue, microenvironment or cancer-context program (cell-extrinsic KEGG or
  Hallmark set); grey, cell-intrinsic KEGG or Hallmark pathway. Lower plots: the per-method demoted versus promoted split across all HGA tissue signatures tested in each task (38 for
  DepMap, 37 for CTRPv2); blue, demoted
  ($\mathrm{NES}<0$); red, promoted ($\mathrm{NES}>0$). In both datasets, linear residualization and within-tissue SHAP demote lineage and tissue-context programs (most HGA tissue signatures have negative NES), whereas IRM shows the weakest separation. In cell lines these tissue-context programs likely reflect lineage-determined expression rather
  than an active microenvironment, so the demotion likely reflects removal of
  the lineage confounder rather than loss of microenvironmental signaling.}
  \label{fig:pathway_delta_dotplot}
\end{figure}
